\documentclass[11pt]{article}
\usepackage[margin=1in]{geometry}
\usepackage{enumitem}
% =============================================================================
% Preambles/header.tex — shared LaTeX/Beamer preamble for this project
%
% Use from a Beamer lecture by prepending:
%     \input{header}
% and compiling with TEXINPUTS=../Preambles:$TEXINPUTS (the /compile-latex
% skill does this automatically).
%
% PALETTE CONTRACT. Color names defined here MUST match the names in
% Quarto/theme-template.scss so Beamer and Quarto renderings use the same
% palette. The scripts/check-palette-sync.sh script enforces this.
%
% To customize: edit the HEX values below AND the matching $variable in
% Quarto/theme-template.scss, then run scripts/check-palette-sync.sh.
% =============================================================================

% -----------------------------------------------------------------------------
% Engine detection + encoding
%
% Our /compile-latex skill uses XeLaTeX, where inputenc is unnecessary and
% can produce encoding warnings. Only load inputenc under pdfTeX.
% -----------------------------------------------------------------------------
\usepackage{iftex}
\ifPDFTeX
  \usepackage[utf8]{inputenc}
\fi

\usepackage{amsmath, amssymb, amsthm}
\usepackage{booktabs}
\usepackage{graphicx}

% -----------------------------------------------------------------------------
% PALETTE — mirrors Quarto/theme-template.scss variable names.
% Load xcolor and define colors BEFORE anything that references them
% (notably hyperref, hypersetup, and the Beamer color assignments).
% -----------------------------------------------------------------------------
\usepackage{xcolor}

\definecolor{primary-blue}{HTML}{012169}      % headings, accents
\definecolor{primary-gold}{HTML}{B9975B}      % emphasis, borders
\definecolor{highlight-yellow}{HTML}{F2A900}  % markers, alerts
\definecolor{light-bg}{HTML}{E8EDF5}          % title / section backgrounds
\definecolor{jet}{HTML}{1A1A1A}               % body text

% Semantic colors (used in diagrams and annotations)
\definecolor{positive}{HTML}{15803D}          % good / observed / identified
\definecolor{negative}{HTML}{B91C1C}          % bad / problematic / bias
\definecolor{neutral}{HTML}{525252}           % reference / context

% Highlight accents (match .hi-* classes in the SCSS)
\definecolor{hi-slate}{HTML}{314F4F}
\definecolor{hi-green}{HTML}{2E8B57}
\definecolor{hi-red}{HTML}{C41E3A}

% Semantic aliases so skill templates and snippets can write
% \textcolor{accent}{...} without hard-coding "primary-blue".
\colorlet{accent}{primary-blue}
\colorlet{accent2}{primary-gold}

% -----------------------------------------------------------------------------
% Hyperref — loaded AFTER xcolor + palette so link colors resolve.
% -----------------------------------------------------------------------------
\usepackage{hyperref}
\hypersetup{colorlinks=true, linkcolor=primary-blue, urlcolor=primary-blue,
            citecolor=primary-blue, pdfborder={0 0 0}}

% -----------------------------------------------------------------------------
% Course code — set once per deck (after \input{header}, before \begin{document})
% via \coursecode{CS401}. Shown in the footer on every slide so a deck's course
% is identifiable without opening the title page. Empty by default.
% -----------------------------------------------------------------------------
\makeatletter
\newcommand{\coursecode}[1]{\gdef\@coursecodetext{#1}}
\gdef\@coursecodetext{}
\makeatother

% -----------------------------------------------------------------------------
% Beamer theme assignments (only apply under Beamer).
%
% We detect Beamer via \@ifclassloaded{beamer}, which is the documented,
% non-internal way. This must be wrapped in \makeatletter/\makeatother
% because \@ifclassloaded contains an @-catcode character.
% -----------------------------------------------------------------------------
\makeatletter
\@ifclassloaded{beamer}{%
  \usefonttheme{professionalfonts}
  \setbeamercolor{structure}{fg=primary-blue}
  \setbeamercolor{frametitle}{fg=primary-blue}
  \setbeamercolor{title}{fg=primary-blue}
  \setbeamercolor{subtitle}{fg=primary-gold}
  \setbeamercolor{author}{fg=primary-blue}
  \setbeamercolor{institute}{fg=neutral}
  \setbeamercolor{date}{fg=primary-gold}
  \setbeamercolor{itemize item}{fg=primary-blue}
  \setbeamercolor{itemize subitem}{fg=primary-gold}
  \setbeamercolor{alerted text}{fg=primary-gold}
  \setbeamercolor{block title}{fg=white, bg=primary-blue}
  \setbeamercolor{block body}{bg=light-bg}
  % Semantic block variants: block=definition/concept (blue), exampleblock=
  % worked example (green/positive), alertblock=key takeaway or caution (gold).
  % Keep to <=2 colored boxes per slide across ALL three variants (INV-7).
  \setbeamercolor{block title example}{fg=white, bg=positive}
  \setbeamercolor{block body example}{bg=light-bg}
  \setbeamercolor{block title alerted}{fg=white, bg=primary-gold}
  \setbeamercolor{block body alerted}{bg=light-bg}

  % Minimal footer: course code (left) + page X / N (right), no navigation clutter
  \setbeamertemplate{navigation symbols}{}
  \setbeamertemplate{footline}{%
    \color{neutral}\footnotesize\hspace{0.7em}\@coursecodetext\hfill
    \insertframenumber/\inserttotalframenumber\hspace{0.7em}\vspace{0.3em}}
}{}
\makeatother

% -----------------------------------------------------------------------------
% TikZ setup — libraries the snippet gallery uses
% -----------------------------------------------------------------------------
\usepackage{tikz}
\usetikzlibrary{arrows.meta, positioning, calc, decorations.pathreplacing,
                fit, shapes.geometric, backgrounds}

% Shared TikZ styles — snippets and hand-written diagrams can pick these up
% via `\begin{tikzpicture}[every node/.style={...}]` or by naming specific
% styles (e.g., `\node[dag-node] (X) at (0,0) {X};`).
\tikzset{
  % Node styles for causal DAGs and similar
  dag-node/.style = {
    draw=primary-blue, fill=light-bg, rounded corners=2pt,
    minimum width=1.2cm, minimum height=0.9cm, align=center, font=\small
  },
  decision-node/.style = {
    draw=primary-gold, fill=white, diamond, aspect=1.8,
    minimum width=1.8cm, minimum height=1.2cm, align=center, font=\small,
    inner sep=1pt
  },
  % Edge styles — observed vs counterfactual
  observed-edge/.style = {-{Latex[length=2.5mm]}, thick, draw=jet},
  counterfactual-edge/.style = {-{Latex[length=2.5mm]}, thick, draw=neutral, dashed},
  confound-edge/.style = {-{Latex[length=2.5mm]}, thick, draw=negative, dashed},
  % Data-point styles for plots
  observed-dot/.style = {fill=primary-blue, draw=primary-blue, circle, inner sep=1.2pt},
  counterfactual-dot/.style = {fill=white, draw=neutral, circle, inner sep=1.2pt}
}

% -----------------------------------------------------------------------------
% Convenience macros
% -----------------------------------------------------------------------------
\newcommand{\muted}[1]{\textcolor{neutral}{#1}}
\newcommand{\key}[1]{\textcolor{primary-gold}{\textbf{#1}}}
\newcommand{\good}[1]{\textcolor{positive}{#1}}
\newcommand{\bad}[1]{\textcolor{negative}{#1}}

% Standout frame macro: full-bleed dark background for section transitions.
% Beamer-only (uses \begin{frame}/\setbeamercolor/\usebeamerfont) -- do not
% call from an article-class file (e.g. Notes/<CODE>/*-notes.tex). \muted,
% \key, \good, \bad, and \coursecode above are plain xcolor/gdef and are
% safe to use from either class; verified when Notes/article-class support
% was added.
% Expand: \transitionslide{Your transition title}
\newcommand{\transitionslide}[1]{%
  \begingroup
  \setbeamercolor{background canvas}{bg=primary-blue}
  \begin{frame}[plain]
    \centering
    \vfill
    {\usebeamerfont{title}\color{white}\Large #1\par}
    \vfill
  \end{frame}
  \endgroup
}

% Section divider frame (Beamer-only), an alternative to \transitionslide with a
% darker two-line style: near-black (jet) background, a small white label line
% above a larger gold title, and a thin gold rule along the bottom edge.
% Expand: \sectiondivider{Section 2}{Pointers}
\newcommand{\sectiondivider}[2]{%
  \begingroup
  \setbeamercolor{background canvas}{bg=jet}
  \begin{frame}[plain]
    \vspace*{3.0cm}
    {\color{white}\ttfamily\large #1\par}
    \vspace{0.5cm}
    {\color{primary-gold}\LARGE #2\par}
    \begin{tikzpicture}[remember picture, overlay]
      \draw[primary-gold, line width=2.5pt]
        ($(current page.south west)+(0,0.1)$) -- ($(current page.south east)+(0,0.1)$);
    \end{tikzpicture}
  \end{frame}
  \endgroup
}

% -----------------------------------------------------------------------------
% Channel watermark — "Concepts That Click" (YouTube).
%
% Article-class files (CompetitiveExam Q/A sets, Notes, Assignments) get a
% light diagonal draft watermark on every page plus a centered footer naming
% the channel. Beamer decks get a subtle TikZ background watermark instead
% (the footline already carries the course code). Centralized here so every
% MCQ PDF is branded without per-file edits.
% -----------------------------------------------------------------------------
\makeatletter
\@ifclassloaded{article}{%
  \usepackage{draftwatermark}
  \SetWatermarkText{Concepts That Click}
  \SetWatermarkScale{1}
  \SetWatermarkFontSize{2cm}
  \SetWatermarkLightness{0.86}
  \SetWatermarkAngle{45}
  \usepackage{fancyhdr}
  \pagestyle{fancy}
  \fancyhf{}
  \fancyfoot[C]{\footnotesize\textcolor{neutral}{Concepts That Click~|~YouTube: \href{https://www.youtube.com/@concepts.thatclick}{@concepts.thatclick}}}
  \renewcommand{\headrulewidth}{0pt}
  \renewcommand{\footrulewidth}{0pt}
  \fancypagestyle{plain}{%
    \fancyhf{}%
    \fancyfoot[C]{\footnotesize\textcolor{neutral}{Concepts That Click~|~YouTube: \href{https://www.youtube.com/@concepts.thatclick}{@concepts.thatclick}}}%
    \renewcommand{\headrulewidth}{0pt}%
    \renewcommand{\footrulewidth}{0pt}%
  }
}{}
\@ifclassloaded{beamer}{%
  \setbeamertemplate{background}{%
    \begin{tikzpicture}[remember picture, overlay]
      \node[rotate=45, opacity=0.055, align=center]
        at (current page.center)
        {\Huge\textcolor{neutral}{Concepts That Click}};
    \end{tikzpicture}%
  }
}{}
\makeatother 
\coursecode{JKSSB}
\renewcommand{\thesection}{66.\arabic{section}}
\newtheorem{example}{Example}[section]
\newenvironment{definitionbox}[1]{%
  \par\noindent\textbf{Definition (#1).}\ \itshape
}{\par\vspace{0.3em}}
\title{Current-Paper Depth Audit and Technical-Post Routing --- Detailed Lecture Notes}
\author{JKSSB Computer Awareness}
\date{\today}

\begin{document}
\maketitle

\section{What this capstone does}

This lesson is a method, not a topic. Earlier lessons taught computer facts;
this one teaches how to keep those facts current and how to study them at the
right depth for the post applied for. The method has two parts. The
\textbf{depth audit} checks notes and practice items against the latest
official paper and answer keys. \textbf{Technical-post routing} directs deeper
material only to the candidates whose notification actually tests it, instead
of generalising Website Operator or Computer Instructor depth to every post.

\begin{definitionbox}{Depth audit}
A repeatable check in which each computer item from the latest official paper
is collected, tagged by depth, compared across the provisional and final
answer keys, screened for version sensitivity, and used to route study time by
post.
\end{definitionbox}

\begin{definitionbox}{Technical-post routing}
The rule that deeper, computational, or engineering-level material is studied
deeply only for the technical posts whose papers contain it, and is read
lightly or skipped by general-post candidates.
\end{definitionbox}

\begin{center}
\begin{tabular}{p{0.26\textwidth}p{0.64\textwidth}}
\toprule
\textbf{Term} & \textbf{Meaning in this lesson} \\
\midrule
Core-General & Stable fundamentals every post tests. \\
Current-General & Stable concept whose current form must be re-checked each cycle. \\
Technical-Post & Deeper material asked only of Website Operator / Computer Instructor. \\
Version-Sensitive & Any claim whose answer depends on a software or protocol version. \\
Tier 1 & JKSSB itself: notifications, papers, provisional and final keys. \\
Final key & The post's final answer key; it settles every dispute for that post. \\
Provisional key & The first key released; subject to change; never memorised as final. \\
\bottomrule
\end{tabular}
\end{center}

\section{Why papers drift}

Three forces make last year's notes stale. First, each recruitment cycle
carries its own notification with its own scope, marks, and penalty rule, so
the boundary of what is asked moves by post. Second, each paper is followed by
a provisional key and then a final key, and answers on disputed or defective
items can change between the two. Third, the subject itself moves: Office
menus and shortcuts shift between versions, and protocol facts (which
transport a version uses, how a mail protocol behaves) depend on the version
named. A note that was correct last year is therefore only a draft until it
has been checked against the newest Tier-1 set.

\section{The five-step revision workflow}

Run these steps once per recruitment cycle, before revising any topic lesson.

\textbf{Step 1: collect the newest Tier-1 set.} Gather the notification and
syllabus scope, the question paper for the target post, and both the
provisional and the final answer key. Record which key is final and which is
still provisional.

\textbf{Step 2: tag every computer item.} Assign each item exactly one depth
tag (definitions in the next section), and add a version flag wherever the
answer could depend on a software or protocol version.

\textbf{Step 3: diff the keys.} Place the provisional key beside the final key
and mark every item where the answer changed, was dropped, or gained an
alternate answer. Changed items are signals of defective or disputed wording,
not new facts to memorise blindly.

\textbf{Step 4: screen for version sensitivity.} For each item about Office,
mail, or protocols, tie the claim to its version using vendor documentation
or the standard. If no version is stated and the answer depends on one, flag
the item as under-specified rather than memorising it.

\textbf{Step 5: route study by post.} Allocate deep study to the tags the
target post actually tests. A Junior Assistant candidate masters the
Core-General tag deeply and reads the other tags lightly; a Computer
Instructor candidate masters the Technical-Post tag as well.

\begin{example}
A candidate holds last year's Office notes and this year's paper. Step 1
collects the new notification, paper, and both keys. Step 2 tags the Office
menu items as Version-Sensitive and the recall items as Core-General. Step 3
finds one answer changed between the provisional and final keys, so that stem
is flagged as disputed wording. Step 4 ties the surviving Office claims to the
prescribed version. Step 5 routes the general-post candidate back to the
Core-General items first.
\end{example}

\section{The four depth tags}

\textbf{Core-General} covers stable fundamentals in the Junior Assistant
Unit IV scope: the bit is the smallest unit of data, RAM is volatile, F5
starts a slide show from the beginning, a formula must start with \texttt{=},
IPv6 is 128-bit. These items are safe to drill for every post.

\textbf{Current-General} covers concepts whose stable core must be refreshed
against current sources: computer virus and latest anti-virus, terms and
abbreviations in current use, and the role of IT in governance (portal and
scheme facts). The concept is permanent; its current instance changes.

\textbf{Technical-Post} covers deeper material found in the Website Operator
and Computer Instructor papers but not in the general scope: cache levels,
encoding details such as UTF-8 variable length, binary conversions with
fractional parts, and statement-evaluation items over engineering concepts.
This tag is studied deeply only by technical-post candidates.

\textbf{Version-Sensitive} is a flag added to any of the above whenever the
answer depends on a version: which tab or menu holds an Office command, the
exact behaviour of a shortcut in a stated version, POP3 download against IMAP
sync, or which transport a protocol version uses. A version-sensitive claim
without a stated version is flagged, not memorised.

\begin{definitionbox}{Tagging test}
Ask in order: (1) does the Junior Assistant scope name this concept --- if
yes, start at Core-General; (2) does the answer change with time, portal, or
version --- if yes, add Current-General or Version-Sensitive; (3) does it need
engineering depth or computation --- if yes, mark Technical-Post.
\end{definitionbox}

\begin{example}
``F5 starts a slide show from the beginning'' is Core-General: stable, in
scope, version-free. ``Which menu holds this Word option'' is
Version-Sensitive: answerable only for a stated Office version. ``Binary
11010101.01 equals 213.25'' is Technical-Post: a computation drill from the
technical papers. ``Latest anti-virus'' is Current-General: the concept is
stable but the instance must be refreshed.
\end{example}

\section{The official source ledger}

Authority flows downward: Tier 1 settles disputes, Tier 2 supplies practice,
Tier 3 explains concepts. The ledger below mirrors the course knowledge base
source registry. No page numbers, paper serials, or key dates are cited here
beyond what the registry records; anything not yet confirmed is marked TBD.

\begin{center}
\begin{tabular}{p{0.14\textwidth}p{0.76\textwidth}}
\toprule
\textbf{ID} & \textbf{Source and what it settles} \\
\midrule
SRC-01 & Notification 08 of 2025 with the official Junior Assistant computer scope (Unit IV) and the 10-mark Computer Applications annexure. Settles what is in scope. \\
SRC-02 & JKSSB question-paper index. Settles which papers exist per year and post. \\
SRC-03 & Junior Assistant 2026 paper, computer block Q61--80. Settles real JA item formats, concepts, and traps. \\
SRC-04 & Website Operator 2026 paper, computer block circa Q61--80. Settles WO item formats and foundational concepts. \\
SRC-05 & Computer Instructor-Operator 2025 paper, 120 items. Settles technical-post depth and the penalty rule. \\
SRC-06 & Junior Assistant 2026 final answer key. Settles disputed JA items. \\
SRC-07 & Website Operator 2026 provisional key. Provisional only; always mark its status. \\
SRC-08 & Computer Instructor/Operator 2025 final key. Settles disputed technical-post items. \\
\midrule
SRC-09/10 & Tier-2 banks (Sanfoundry and standard compilations). Practice composition and distractor mining only; never authoritative; label adapted items accordingly. \\
SRC-11/12/13 & Tier-3 references (portal pages, vendor docs, textbooks). Concept exposition only; vendor claims tied to the prescribed Office version. \\
\bottomrule
\end{tabular}
\end{center}

The single most important rule of the ledger is precedence: a Tier-1 final
key supersedes every bank, tutorial, and textbook on a disputed item. A lower
tier is never used to overrule a higher one.

\section{How to check answer-key changes}

Take the provisional key and the final key for the same post and compare them
item by item. For each computer item, record one of four outcomes: unchanged,
changed, dropped, or given an alternate answer. For every item that is not
unchanged, go back to the stem and quote its option wording verbatim, then
flag the defect (ambiguous wording, two defensible answers, version-free menu
claim) and teach the stable concept underneath instead of memorising the new
letter. Keep this diff with the notes as the cycle's audit trail, and never
carry a provisional answer into a mock once the final key has been published.

\begin{center}
\begin{tabular}{p{0.22\textwidth}p{0.68\textwidth}}
\toprule
\textbf{Key-diff outcome} & \textbf{What to record and do} \\
\midrule
Unchanged & No action; the item confirms the notes as written. \\
Changed & Quote the stem verbatim; flag the wording defect; teach the stable concept; learn the final letter. \\
Dropped & Quote the stem verbatim; note the topic as dispute-prone; do not drill the dropped letter. \\
Alternate answer allowed & Note both acceptances; teach the concept so either option is derivable. \\
\bottomrule
\end{tabular}
\end{center}

Two cautions belong here. First, a changed answer never means ``the concept
changed''; it means the wording admitted two readings and the board has now
picked one. Students who memorise only the new letter without the flagged
wording will fall for the same trap when it returns in paraphrase. Second, a
provisional key is never quoted as authority in a dispute with a fellow
candidate; only the final key, with its post and status named, settles the
argument.

\section{Why technical-post depth must not be generalised}

The Junior Assistant paper is dominated by recall and classification at the
intro-to-core level, while the Computer Instructor paper adds computation and
engineering depth. Generalising the deeper block to every candidate has two
costs: the general-post candidate wastes time on material the notification
never tests and absorbs L4/L5 distractors, while the technical-post candidate
who studies only the general block arrives under-prepared. The correction is
the routing rule: tag each item first, then allocate study time by tag. Binary
conversions with fractional parts, encoding internals, and cache-level
internals are Technical-Post; they are required for Computer Instructor,
optional background for Junior Assistant. Foundational recall (bit as smallest
unit, ISP expansion, \texttt{@} separating user name and domain name) is
Core-General; it is required for every post.

\section{Version sensitivity: Office and protocols}

Office items split into two kinds. Stable Office facts (Mail Merge combines
letters with an address database; a formula must start with \texttt{=}; a
Select query retrieves data on a condition) are Core-General and safe to
drill. Menu-path, tab-location, and shortcut-detail items are
Version-Sensitive: they are trustworthy only when tied to the prescribed
Office version through vendor documentation, and an item that names no version
but depends on one is flagged as under-specified.

The table below works the tagging test through on items anchored in the
course's real-PYQ evidence registry, so the routing decision is visible in one
place. Each row gives the item, its tag, and the reason in one line.

\begin{center}
\begin{tabular}{p{0.40\textwidth}p{0.20\textwidth}p{0.30\textwidth}}
\toprule
\textbf{Item} & \textbf{Tag} & \textbf{Why this tag} \\
\midrule
Bit is the smallest unit of data & Core-General & Stable recall, in every post's scope \\
F5 starts a slide show & Core-General & Stable, version-free behaviour \\
ISP expands to Internet Service Provider & Core-General & Exact abbreviation recall \\
Latest anti-virus & Current-General & Concept stable, instance refreshable \\
Role of IT in governance & Current-General & Portal and scheme facts change \\
Binary 11010101.01 to decimal & Technical-Post & Computation drill, technical papers \\
UTF-8 variable-length, ASCII-compatible & Technical-Post & Encoding internals, technical papers \\
Which tab holds this Word command (no version) & Version-Sensitive & Answer moves between Office versions \\
HTTP version transport claim & Version-Sensitive & Depends on the protocol version named \\
\bottomrule
\end{tabular}
\end{center}

A Junior Assistant candidate drills the first three rows deeply, refreshes the
next two against current sources, reads the technical rows lightly, and
refuses to memorise the version-sensitive rows until a version is stated. A
Computer Instructor candidate adds deep drill on the technical rows. Neither
candidate generalises the other's depth.

Protocol items behave the same way. The stable splits do not move with
versions: IPv6 is 128-bit with a 40-byte header, POP3 downloads to the local
machine while IMAP syncs, HTTP/3 uses QUIC over UDP rather than TCP, a
digital signature provides authenticity and integrity while encryption
provides confidentiality, and BCC hides recipients from each other. What is
version-sensitive is any claim about a specific protocol version's transport
or behaviour; those claims are settled by the final key together with the
standard, never by a bank alone.

\section{Exam retrieval and common confusions}

Five audit-stage confusions cause most errors on this kind of material.
\begin{enumerate}
  \item A \textbf{provisional key} is treated as final. The correction is to
    mark its status and re-check every disputed item against the final key.
  \item A \textbf{generated practice item} is treated as a real PYQ. The
    correction is the provenance label: generated items carry
    \textit{[Original]}; only verbatim official items carry a PYQ label.
  \item \textbf{Technical-post depth} is studied as if every post tested it.
    The correction is the tag: Technical-Post material is deep study only for
    the technical posts.
  \item A \textbf{defective option} is memorised as fact. The correction is to
    quote it verbatim, flag the defect, and teach the stable concept.
  \item A \textbf{Tier-2 bank} is used to overrule a Tier-1 key. The
    correction is precedence: the final key always wins.
\end{enumerate}

\begin{example}
An item states, ``POP3 keeps mail synced on the server, so it is the best
choice for multi-device sync.'' The routing check tags the concept as
Version-Sensitive and the claim as false under the stable split: POP3
downloads to the local machine while IMAP syncs. The bank that suggested the
stem does not overrule the standard or the final key, and the defective
wording is flagged rather than memorised.
\end{example}

\subsection*{Exercises}
\begin{enumerate}[label=\arabic*.]
  \item State the five steps of the revision workflow in order.
  \item Define the four depth tags and give one example of each.
  \item List the Tier-1 ledger entries and state what each one settles.
  \item Explain how to check provisional-versus-final key changes and what to do with a changed item.
  \item Explain why Website Operator / Computer Instructor depth must not be generalised, with one example.
  \item Distinguish a stable Office fact from a version-sensitive Office claim, with one example of each.
\end{enumerate}

\subsection*{Solutions}
\begingroup\small
\begin{enumerate}[label=\arabic*.]
  \item Collect the newest Tier-1 set; tag every computer item; diff the provisional against the final key; screen for version sensitivity; route study by post.
  \item Core-General: stable fundamentals every post tests (F5 starts a slide show). Current-General: stable concept with a refreshable instance (latest anti-virus). Technical-Post: deeper material for technical posts (binary conversion drills). Version-Sensitive: answer depends on a version (menu-path or protocol-transport claims).
  \item SRC-01 notification and scope; SRC-02 paper index; SRC-03 JA 2026 paper; SRC-04 WO 2026 paper; SRC-05 CI 2025 paper; SRC-06 JA final key; SRC-07 WO provisional key (provisional only); SRC-08 CI final key. Papers settle formats and concepts; final keys settle disputes.
  \item Compare the two keys item by item; mark changed, dropped, or alternate-answer items; quote the stem verbatim, flag the defect, teach the stable concept; never carry a provisional answer into a mock once the final key exists.
  \item Because the general papers test recall and classification while the technical papers add computation and depth; generalising wastes general candidates' time and under-prepares no one correctly. Example: fractional binary conversion is required for Computer Instructor but optional for Junior Assistant.
  \item Stable: Mail Merge combines letters with a database, or formulas start with \texttt{=}. Version-sensitive: which tab holds a command, or the exact shortcut detail in a stated version --- trustworthy only when tied to the prescribed version.
\end{enumerate}
\endgroup
\end{document}
